\noindent\textbf{1. How does the feeding law change coexistence?}\par
We kept the coursework's weak Allee prey growth and changed only the feeding law. All five models preserve positive populations and keep them bounded. In the two prey-dependent models (P1 and P2), positive coexistence fixes prey density independently of carrying capacity. The original foraging-arena response (P0) can support two attracting coexistence states. With ratio-dependent feeding (P4), enrichment lowers the Jacobian trace at every positive coexistence state with positive determinant for the coursework growth law. For a positive growth factor $g$ with $g'>0$, the same decrease holds at a given prey density for every admissible ratio-dependent intake if and only if $(1/g)''\geq0$ there.

\medskip
\noindent\textbf{2. How do prey growth and predator competition interact?}\par
We kept the feeding law fixed and paired three prey growth laws with and without direct predator competition. Logistic growth allows at most one positive equilibrium. Weak and strong Allee growth can each support three. When the three growth laws are calibrated to the same coexistence state at zero competition, positive competition orders predator density along their common branches with positive determinant. Predator density is highest for strong Allee growth, intermediate for weak Allee growth, and lowest for logistic growth. This ranks equilibrium levels even when predator density falls in all three models. In 48 disturbance simulations at the study baseline, 45 recovered. All three failures followed prey removal under strong Allee growth.

\medskip
\noindent\textbf{3. What transfers when the same equation is fitted to other systems?}\par
We kept the ecological equation and fitted its coefficients to training conditions in tumor--T-cell and bacteria--phage datasets. The equation had lower mean held-out error than the CARRGO model on the PBT138 cell line, while the infection-stage model had lower four-hour error in 11 of 12 $\lambda$-phage conditions. In 12 held leukemia conditions, fitting targets alone predicted none of the observed early T-cell increases. Fitting both populations predicted all 12, but mean target root-mean-square error rose from 81,634 to 197,580 cells. The measured trade-off concerns predictions for held conditions from this dataset.
